\documentclass[11pt]{article}
\usepackage[margin=1in]{geometry}
\usepackage[T1]{fontenc}
\usepackage{lmodern,booktabs,graphicx}
\usepackage[hidelinks]{hyperref}
\usepackage{xurl}
\setlength{\parindent}{0pt}
\setlength{\parskip}{6pt}
\setlength{\emergencystretch}{3em}
\title{Native Perception and Estimation After Hiding Parked Vehicles}
\author{Pan Dynamics Collision Avoidance Research}
\date{October 2, 2026}
\begin{document}
\maketitle

\section{Result}
A fresh CARLA capture with parked vehicles hidden completed causal sensor
preparation, actual YOLO26x car detection, LiDAR rectangle localization and
unchanged NRMM estimation using native RGB, LiDAR, GNSS and IMU measurements.
All three predeclared cases used confidence association; persistent association
was not enabled. The three replays produced 900 finite estimator publications
with matching frame IDs and timestamps. Cases passing all existing smoke
accuracy gates: 3/3. Gains and thresholds were not retuned
after scoring this capture.

\begin{center}
\small
\begin{tabular}{@{}lrrrrl@{}}
\toprule
Case & Raw pos. & Est. pos. & Est. vel. & Valid & Gates \\
 & (m) & (m) & (m/s) & frames & \\
\midrule
Original filters & 0.3695 & 0.4641 & 1.2201 & 300/300 & Pass \\
Improved filters & 0.1629 & 0.1831 & 0.4195 & 300/300 & Pass \\
Improved + blackout & 0.1555 & 0.2054 & 0.4804 & 285/300 & Pass \\
\bottomrule
\end{tabular}
\end{center}

Errors are vector RMSE over evaluation times 2.00--5.98 s, excluding the
first two seconds of estimator acquisition. Raw-perception RMSE includes
only frames with an available measurement; estimator RMSE includes missing
measurement intervals. Target gates remain 0.5 m position and 1.5 m/s velocity.
Raw-position gates remain 0.5 m RMSE and 1.0 m p95, with at least 95\%
availability outside the deliberately blanked interval. The original ego
position, velocity, body-velocity and yaw gates are retained in
\path{protocol.json}. All three cases share exactly the same ego estimates.

For comparison, the original parked-car capture with confidence association
and original filters had raw/estimated-position/velocity RMSE of
1.2456 m, 1.2406 m and 3.6722 m/s. The previous improved-filter
capture with parked cars and persistent association had 0.2851 m position
and 0.7453 m/s velocity RMSE. The repeated route and seeds do not establish
independent scenario generalization. Over all 500 native steps, the old and
fresh ego positions and velocities were identical. Maximum absolute target
position and velocity differences were 0.00007439 m and 0.00352669
m/s, respectively. This comparison is scoring-only.

\section{Fresh capture and fixed prediction settings}
An owned CARLA 0.10.0 server ran Town10HD\_Opt on port 2011. The native
capture driver loaded \path{config/carlaPerceptionScene.json} and disabled
143 static vehicle mesh components before spawning the Lincoln MKZ ego and
Dodge Charger target. These are component counts, not distinct vehicle counts.
The map-selected straight used spawn index 105, a 15 m initial gap and
11 m/s nominal waypoint-following speeds. Fifty stationary calibration and
150 driving warm-up samples preceded 300 evaluation samples at 0.02 s.
All 500 RGB/LiDAR/GNSS/IMU groups had matching sensor IDs and consecutive
simulation frame IDs.

The front camera was 640-by-640 with a 100-degree field of view. LiDAR had
64 channels, 1,280,000 points/s, 50 revolutions/s, 80 m range and 0.02 m
range-noise standard deviation. Native GNSS used approximately 0.01 m
position noise; IMU acceleration/gyro standard deviations were 0.01 m/s$^2$
and 0.0005 rad/s. LiDAR/IMU seed was 20261002; GNSS seed was 20261003.
Capture took 51.522 s. Cleanup restored the hidden environment
objects and asynchronous settings, left zero vehicle/sensor actors, and
stopped only the verified owned PID with SIGINT. Exit code 130, absent PID
and closed port were checked. Shared port 2000 was not connected to, ticked,
modified or signaled.

Preparation used only native GNSS/IMU, fixed georeference and sensor
extrinsics. The existing causal 0.4 s endpoint GNSS derivative, stationary
tilt initialization, gyro propagation, native compass yaw and gravity removal
were unchanged. An audit hook prohibited preparation from opening the scoring
truth file. Prediction manifests contain no actor-truth state. All perception
summaries recorded zero truth comparisons; all six prediction-file hashes
were written before the separate scorer opened truth.

Every case used YOLO26x, confidence 0.18, NMS IoU 0.45, a front camera,
causal rectangle fitting and declared 5.0-by-1.9 m target dimensions.
The original filters used minimum depth 1.0 m and minimum object height
0.3 m. The improved filters used 0.1 m and 0.0 m, respectively, with the
same 1.2 m cluster radius. These two filter changes were selected in the
preceding experiment, before this new run. No model weights, estimator source,
estimator configuration or native integration binary changed during testing.

\section{Missing measurements, selection and limits}
The third case changed only the derived camera paths for frames 150--164
(3.00--3.28 s) to an existing verified black image. The native recordings,
LiDAR and ego inputs remained intact. The actual detector produced
15 missing measurements in that interval, and reacquired at frame
165. Estimation remained finite; maximum target-position error during
the interruption was 0.4783 m. The first 150 perception outputs and
estimator states were exactly identical to the improved-filter normal case
(excluding timing). All 14 existing Python association/causality tests passed.

Scoring-only projection of the target center found selected-box center
exclusion counts of 0, 0, 0 in the three cases. Center inclusion is a
selection diagnostic, not a complete identity or bounding-box annotation.
Visual inspection of new RGB frame 237 showed the target retained and
roadside cars removed.

This is a functional and numerical accuracy replay, not a certified error
bound or real-time closed-loop result. Raw relative-position errors exceeded
the configured 0.12 m deterministic measurement bound on 294, 230, 215
frames, respectively. The observed target-speed range was 4.757--11.141
m/s; the configured target domain starts at 10 m/s. Availability of a runtime
bound flag does not verify these measurement/domain premises. Native compass
yaw also retains the preceding experiment's ideal-heading limitation.
The runtime flagged the estimated operating domain invalid on 179, 127, 127
frames, respectively, despite making position-bound values available.

The three perception processes, including model loading, native-file reads
and fitting, took 49.373, 99.994, 92.256 s for 300 frames each. This offline
throughput does not demonstrate 50 Hz operation. The MATLAB batch completed
three unchanged-estimator replays in 17.760 s. Per-sample estimator
timing is reported separately in the score and runtime summaries.

\section{Artifacts and reproduction}
The experiment plan, exact capture/pipeline commands, compact scores,
configuration, source and model hashes, causal-prefix audit and cleanup checks
are in \path{report/PARKED_FREE_PERCEPTION_20261002/}.
Original native data and experiment drivers are under
\path{simulation_output/parked_free_perception_20261002/}; these local CARLA
artifacts remain untracked. The external archive preserves export copies,
the project commit record and original-source hashes. Generated plots/PDFs
are local and archived outputs, not tracked data or binaries.

\end{document}
